\chapter[Philosophy Is Not Memorising Answers]{Philosophy Is Not Memorising Philosophers' Answers}
\section{An Answer Sheet without Boxes}
Pick up something utterly familiar: an answer sheet.

Barely larger than your palm, it carries neat arrays of pale blue or pink boxes, with four little ovals for each question: A, B, C and D. You fill your chosen oval with a 2B pencil, darkly and completely, without crossing its boundaries. A marking machine scans a hundred questions in a second and declares right or wrong, without negotiation.

Study its construction. It is really a declaration of a worldview: every question has exactly one correct option; options have clear boundaries; assessing correctness requires neither knowing who you are nor why you think as you do. Most importantly, every worthwhile question has already been asked. Your job is merely to answer.

Imagine someone leaving every oval blank and writing on the reverse: ``The word `fair' has two meanings in this question. Without clarification, both A and C are correct.'' What will the machine do? Reject the sheet as invalid.

This chapter concerns the space without boxes. Here questions themselves can be examined, dismantled and doubted. Answers are not selected from four options: you build them step by step, then let others try to take them apart. This space has an ancient name: philosophy.

Many imagine a philosophy lesson as memorising a comparison chart: what Socrates, Plato, Zhuangzi and Kant said, then matching names and answers on a sheet. That is philosophy's least important, perhaps even dispensable, part. What really matters will not fit any box.

\section{An Athenian Court in 399 BCE}
Come into a scene.

It is 399 BCE in Athens. The day's examination room is a courtroom, but unlike an answer-sheet examination, it has neither an examiner nor standard answers. Instead, roughly five hundred jurors have been chosen by lot from the citizenry: not legal experts, but potters, farmers, boatmen and small traders with calloused hands. Sun beats on the open-air meeting place and the crowd hums. A water clock measures speeches: water drips steadily between vessels, and when it runs out, speaking time ends.

The accused is a seventy-year-old man with a snub nose, a large belly and bare feet: Socrates. There are three charges: disrespecting the city's gods, introducing new gods and corrupting the young. The first two sound religious; the third is crucial. For decades, he has stopped young people and distinguished citizens in the marketplace, asking embarrassing questions. You claim to know courage: define it. You intend to govern: what is justice? Within a few exchanges, respondents often contradict themselves and spectators burst out laughing. Many targets are prominent citizens.

Athenian procedure lets a convicted defendant propose an alternative penalty for jurors to choose against the prosecutor's demand. The prosecution wants death. How does Socrates use his chance? He declares that his lifelong questioning has served the city as an unpaid gadfly, stinging people awake. His proper treatment, he suggests, is free maintenance in the civic hall, like athletes who bring the city glory.

Imagine those five hundred Athenians' faces. The water keeps dripping. The final vote: death.

Plato records his courtroom stance in the Apology, broadly: an unexamined life is not worth living. Weeks later, Socrates accepts a cup of hemlock in prison and drinks calmly, still discussing with his pupils before death. He retracts none of his questions.

Remember him. We shall return repeatedly to the great paradox he embodies, our entrance into this chapter.

\section{If Books Have the Answers, Why Think?}
Set out the question without hurrying to answer.

Here is the paradox. Socrates, revered throughout Western philosophy as an originator, never wrote a word. No books, papers or standard answers. We know him through records by Plato and others. He left not answers but a way of questioning.

Consider his work. He claimed to know only that he knew nothing. He visited reputedly wise politicians, poets and craftspeople not to seek answers, but to show that they too could not explain themselves clearly. If philosophy's value lies in its output of answers, Socrates was its least successful producer: output zero.

So here is the question:

\textbf{If philosophy has value, where does that value lie?}

Perhaps in philosophers' answers: the world as a shadow of Forms, Plato; the equality of all things, Zhuangzi; ``I think, therefore I am'', Descartes; human beings as ends, Kant. Then the best learning strategy is memorisation: tabulate, learn, fill boxes, score full marks, finish. Teach philosophy like the periodic table.

Or perhaps in the process of reaching answers. Answers are footprints left by other thinkers; philosophy teaches walking: clarifying vague words, testing claims against counterexamples, constructing sound arguments and dismantling others'. Then memorisation can be harmful as well as useless, letting you imagine you can walk because you own a roomful of footprints.

This is not idle academic talk. It determines how philosophy, history, language and all humanities should be taught and examined, and what your time studying them buys.

Before proceeding, choose a side. Who seems more to be learning philosophy: a student who knows philosophers' answers by heart, or one who remembers few but always asks what justifies a claim?

\section{Both Sides Have Reasons}
First present the memorisers' case fully and seriously. Popular discourse has made memorisation synonymous with exam-driven education, unfairly.

They have at least three good reasons. First, knowledge accumulates; each generation need not start from nothing. Newton credited giants' shoulders. Nobody requires you to re-prove Pythagoras before using his theorem. Earlier conclusions form a bridge; demolishing it so everyone must swim is wasteful, not profound. Second, knowing earlier thought prevents repeating elementary errors. Your insight after three days' reflection may have appeared two millennia ago and been refuted four centuries afterwards. Reading earlier answers locates your idea in history's queue. Third, examinations need fairness. Without definite answers, what determines marks: the examiner's mood? Rigid memorisation at least supplies a common ruler. Disadvantaged children can score by hard learning. Remove standard answers and advantage may shift to whoever has a tutor skilled at sounding profound. That would be unfair.

These arguments are substantial. Many of humanity's oldest scholarly traditions operated this way. Chinese classical scholars first mastered and memorised texts and authoritative commentaries, such as Zhu Xi's Collected Commentaries on the Four Books, before examinations and their own interpretations. Swallow answers before digesting them: a standard route for millennia.

Now fully present the case for doing philosophy.

First, philosophical answers differ from mathematical ones. Memorised Pythagoras works whether you understand its proof or not. What can memorising the unexamined-life maxim do? It cannot automatically improve your life; you must examine it yourself. Philosophy's instructions say: take personally; substitution ineffective. Second, memorised answers impersonate thinking. Someone says ``Heidegger said'', listeners nod, and the speaker feels understanding. Real understanding survives three follow-up questions. Quotable maxims are thought's artificial sweetener: sweet, filling, without nourishment. Third, philosophical skills transfer. You probably will not discuss Plato daily, but will encounter vague concepts, dubious reasoning and attractively packaged nonsense. Learning to dismantle these lasts longer than any remembered conclusion.

Other traditions reveal a worldwide tension between process and memorisation. Ancient India's Nalanda, among the greatest scholarly centres, was famous for public debate. Monks had to answer opponents immediately before an audience; failing meant losing, with memorisation no rescue. Medieval European universities required public disputations before graduation, professors and spectators taking turns to object. Degrees went to those who withstood challenge, not those remembering most. Chinese Chan Buddhism went further: a teacher might answer with an apparent non sequitur, a shout or a blow, breaking expectations of standard answers and forcing personal thought.

The real disagreement is not whether answers help: they are maps. It is whether philosophy produces maps or the ability to read and draw them. A memorised map leads where others have gone. The skill can take you where maps remain undrawn.

\section[Thinking Bridge: Four Tools]{Thinking Bridge: The Philosopher's Four Tools}
If philosophy teaches walking rather than memorising footprints, can walking be taught in parts? Yes. Over two millennia, philosophers have developed four everyday tools. Nothing mysterious: take them today.

\textbf{Tool one: clarify concepts.}

Most argumentative impasses arise when one word names different things. Gongsun Long, a Warring States disputant, famously argued that a white horse is not a horse. Apparently quibbling, he was fixing the uses of horse, a category; white, a colour; and white horse, their combination. Different ranges should not be confused. You may reject his conclusion while valuing the move: \textbf{before arguing, ask what the other person means by the word}. Practise by requesting one example inside its scope and one outside. Before debating whether gaming wastes time, define waste. Seven arguments in ten dissolve here: the parties were discussing different things.

\textbf{Tool two: find counterexamples.}

One counterexample kills a universal claim. Europeans were confident for over a millennium that all swans were white, until black swans appeared in Australia. One sufficed. Counterexamples need not speak louder; they need only exist. Against ``Effort guarantees success'', find an unsuccessful hard worker and revise to effort improving the probability. Against ``The internet makes people stupid'', find someone learning a skill online and narrow the claim. The discipline is not winning, but improving accuracy. Whoever rejects revision after a counterexample and accuses the other person of quibbling has truly lost.

\textbf{Tool three: construct arguments.}

An argument supports a conclusion with reasons or premises. Check two things separately: are the premises true, and would the conclusion follow even if all were true? Consider: all cats have tails; my dog has a tail; therefore my dog is a cat. Both premises are true, yet the inference fails. Truth and validity differ. Civilisations independently built testing devices: Greek syllogisms and Indian formats of claim, reason and example. For instance, the hill has fire, because it has smoke, as a smoking stove has fire. Different machinery serves one purpose: conclusions must not stand unsupported.

\textbf{Tool four: thought experiments.}

Especially characteristic of philosophy, these build simplified mental scenes to test claims. No laboratory, no financial cost, considerable power. Consider three famous cases:

\begin{itemize}
\item Galileo may have been the physicist most adept at them. Against heavier bodies falling faster, imagine tying a small stone to a large one. The smaller should retard the larger, slowing the pair; yet their greater combined weight should accelerate it. One claim yields opposite conclusions and therefore fails. No tower climb is required.
\item Plutarch records Theseus's ship, its planks replaced one by one until none remain. Is it the same ship? This targets personal identity too: your cells are replaced in successive batches.
\item Philippa Foot proposed the trolley problem: an uncontrolled trolley approaches five people. You can divert it, but one person occupies the other track. Should you? It presses rival intuitions about consequences and the nature of actions.
\end{itemize}
Thought experiments expose vague claims, but have a limitation. Daniel Dennett called such devices intuition pumps. What they pump is intuition, not an instrument reading. Recent experimental philosophy surveys ordinary intuitions, finding answers shifting with wording and differing across cultures. A thought experiment is a probe, not a judge: it reveals problems rather than pronouncing verdicts. Store that warning for later.

Four tools: clarify concepts, find counterexamples, construct arguments and conduct thought experiments. Together they specify what doing philosophy does.

\section[Debating above the Hao]{Above the Hao: A Debate Two Millennia Ago}
Read a primary source to see these tools at work. Zhuangzi's Autumn Floods records a famous encounter between Zhuangzi and his friend Huizi:

\begin{quote}
Zhuangzi and Huizi strolled above the Hao. Zhuangzi said, ``The minnows swim freely at ease. That is fishes' happiness.'' Huizi said, ``You are not a fish. How do you know their happiness?'' Zhuangzi replied, ``You are not me. How do you know I do not know their happiness?'' Huizi said, ``I am not you, so admittedly I do not know you. But you certainly are not a fish, so your not knowing their happiness is established!'' Zhuangzi said, ``Return to the beginning. When you asked how I knew their happiness, you already knew I knew and asked me. I know it here above the Hao.''
\end{quote}
In brief, the friends reach the bridge. Zhuangzi calls the leisurely fish happy. Huizi questions how a non-fish could know. Zhuangzi asks how a non-Zhuangzi could know his ignorance. Huizi accepts not knowing Zhuangzi and applies the same rule: neither can Zhuangzi know fish. His logical circle closes. Zhuangzi returns to the opening, reading the question as where he knows fishes' happiness, which presupposes knowing. His answer: here on the bridge.

Many language textbooks include this passage, usually declaring Zhuangzi the winner through a higher outlook. Reapply the four tools and a readily misjudged counterexample emerges.

Huizi's opening is a counterexample-style challenge to knowledge's limits across species, a genuine question still present in animal pain and AI feeling research. Zhuangzi's reply cleverly turns the opponent's weapon: claiming to know another's ignorance violates the rule that difference prevents knowledge. But Huizi's third sentence tightens his rule. He concedes his own ignorance; Zhuangzi enjoys no superior position and likewise cannot know fish. Structurally, Huizi now has the stronger argument.

How does Zhuangzi finish? By changing the sense of an zhi: Huizi means how could you know, challenging entitlement; Zhuangzi reads where do you know, requesting a location, and answers on the bridge. This is equivocation, exposed by clarification. Graceful, poetic and beautiful, but under argumentative rules it escapes through ambiguity.

The scoreboard therefore reads: literary applause for Zhuangzi, tighter reasoning for Huizi. Most readers over two millennia, including those who memorised the lesson, learned the result backwards. Memorise conclusions without examining processes and even the winner may be wrong. Still, give Zhuangzi his full case: perhaps he is not refusing defeat, but regards knowing fishes' happiness as felt communion rather than inference. That different philosophical position deserves attention, but needs its own argument.

\section{Why Do Philosophers Keep Arguing?}
A fact now requires explanation.

Many disputes have ended. Observation and telescopes settled whether Earth orbits the Sun; experiments settled whether heavier bodies fall faster; statistics settled bloodletting's medical value. Yet justice, knowledge, mind and whether machines think remain contested after two millennia among brilliant participants. Why can no single experiment commonly end a philosophical dispute?

At least three explanations contend, each with reasons and weaknesses, not interchangeable.

\textbf{First: philosophical problems are pseudo-problems created by confused language.}

Twentieth-century logical positivists held much metaphysical dispute arose from imprecise words and unverifiable sentences mistaken for real questions. Wittgenstein's more famous image likened philosophical problems to language falling ill: clarify them until they dissolve rather than answer them. Many disputes genuinely vanish under clarification. Much argument about numbers' existence, for instance, comes from differing mathematical and everyday uses of existence. Yet some disagreements survive. Clarify fairness endlessly and Rawls and Nozick remain opposed while understanding one another. Diagnosing all disagreement as linguistic disease does not cure every patient.

\textbf{Second: philosophical questions concern frameworks, which experiments cannot adjudicate.}

Experiments operate inside conceptual frameworks; philosophy contests the frameworks themselves. Experiments presuppose evidence, personal identity and causation. Using them to decide those concepts is circular, like candidates setting their own examinations. Measuring planks cannot settle Theseus's ship because the question concerns sameness. Voting cannot settle a trolley case because it counts intuitions, not rightness. Yet overstated, this becomes anything goes without an umpire. That is false. Arguments still contain gaps, concepts contradict themselves and counterexamples defeat claims. No final judge does not mean no court.

\textbf{Third: philosophy progresses, but solved areas gain new names and independence.}

Consider its household's divisions. Physics was natural philosophy, as Newton's Mathematical Principles of Natural Philosophy attests. Logic, psychology and linguistics successively established independent homes. Philosophy's ripened fruit is picked and given away, leaving only unripe, difficult questions in its garden, making harvests seem absent. Optimistic and historically grounded, this nevertheless has limits. Value, meaning and how to live may never mature into separate fields because they involve questioners themselves. Explaining everything as not yet independent may merely postpone disappointment.

Notice two things. These are not peacefully compatible claims that everyone is right. The first calls most problems misunderstandings; the third says they are progressively solved. Both cannot be wholly right. Comparison identifies what each explains and where it stalls, not compromise for its own sake. Second, explanations have standards: coverage, consistency and testable implications. These extend the argument-construction tool.

\section{Further Challenge: Calibrating Doubt}
Now our weightiest theoretical question: how far should doubt go?

Begin with its strongest version. Descartes emptied out every belief, checking each and placing anything doubtful in a suspect pile. Senses deceive: chopsticks look bent in water. Surely writing beside the fire is certain? Dreams feel that way too. What about two plus three equalling five? Imagine an omnipotent demon interfering with every calculation. This is a thought experiment conducted in an armchair, costing nothing and ranging over all knowledge.

Peeling away layers, he found something irreducible. Even deception requires an I to be deceived; doubting demonstrates a doubter. His Discourse on Method states: ``I think, therefore I am.''

Often mistaken for someone committed to doubting everything, Descartes aimed instead to \textbf{find foundations}, not inhabit ruins. Doubt was his bulldozer, not his house. That distinction matters: \textbf{doubt is not a fixed position, but an adjustable dial}.

Imagine belief as radio volume. Strong evidence: turn up. Weak evidence: turn down. New evidence: adjust again. Insufficient evidence and suspended judgement form a legitimate, stable setting, neither failure nor compromise but an honest report. Proportioning belief to evidence is doubt's proper use.

Pause at a counterexample often misjudged. Distrusting everything may seem the clearest protection against deception. In practice it can do the opposite. Equal distrust removes \textbf{discernment}, not credulity. Official accounts and revelations from unidentified, unverifiable accounts become level, inviting emotional choice: believe whatever feels satisfying or distinctive. Conspiracy theorists almost all begin with questioning everything and end staking everything on the least testable accounts. Total doubt dismantles remaining defences rather than curing gullibility. Sophistication means distributing trust by evidence, not disbelieving all.

Doubt also has a logical self-destruct point. What tools permit doubting everything? Thought needs language, logic and the fact of present thinking. Doubt those away too and doubt cannot begin. Absolute doubt cannot support even its own declaration. Healthy doubt is therefore local and anchored: provisionally trust some things to test others, then exchange roles. Like repairing a ship at sea, replace planks individually rather than dismantling the whole vessel at once.

Retrieve our warning: thought experiments are probes, not judges. The brain in a vat modernises Descartes's demon, imagining a brain supplied with computer-generated experience. It can establish that the world might be unreal, but cannot decide whether you should consequently neglect living. A probe reveals cracks; their size, importance and practical implications require calibrated judgement. Philosophy honestly neither promises impregnable foundations nor urges jumping from a building because foundations have cracks.

\section{An Age of Too Many Answers}
This ancient question appears around you in three forms.

First, school. Essays have universal stock examples, comprehension questions templates, even personal opinions lists of scoring points. Genuine questions lose marks by straying, wasting time and lacking boxes. Gradually, thinking too much becomes examination's dearest luxury and memorising without asking its best-value strategy. Our opening answer sheet is mass-produced daily.

Second, the internet. Comment sections wield quotations like hammers: a famous person said this, so stop arguing. Philosophers' names end rather than begin thought, raising a banner of celebrity backing. Meanwhile, changing one's mind counts as defeat. Statements differing from three years earlier invite public exposure. Yet evidence-calibrated belief requires readiness to change. Praising never changing welds doubt's dial in place.

Third, newest and sharpest, something beside you can recite every philosopher's answers. Chatbots fluently explain Plato's cave, Zhuangzi's butterfly and Kant's categorical imperative, faster and more comprehensively than any diligent student, without sleep. Knowing answers has, for the first time, become almost costless.

This puts our disagreement before everyone. If philosophy's value is answers, machines have wholesaled it and humanities education can close. If its value lies in clarification, argument-checking and calibrated doubt, these skills become more valuable. Machine answers share a characteristic: equal fluency. Correct and incorrect claims, supported judgements and fabricated quotations, uncertainty and certainty arrive equally articulately. Amid excess answers and packaged certainty, scarcity reverses. Who asks good questions, detects flawed questions and pauses fluent nonsense? These abilities have no answer-sheet boxes and evade marking machines.

Good questions outweighing hasty answers is therefore no consolation, but a factual claim. Mass-producible answers depreciate; good questions appreciate because they resist mass production. They require genuinely being stuck, knowing where understanding fails and patiently articulating that failure. What Socrates did in the marketplace remains difficult for machines: they hurry to answer.

\section{Your Turn: Which Would You Sign?}
Return to the court in 399 BCE. Socrates had a choice: promising to stop questioning and live quietly might well have spared him. He refused, drank hemlock and left many questions with a bibliography containing no answers: zero books.

Over two millennia later, his questions survive in disputes about happy fish, ships' identity and trolley switches, while few names of the certain jurors remain. Answers seem mortal, questions durable. But hear the other side. Jurors may have wanted more than stupid tranquillity. Cities require decisions, judgements and rules, not endless questions. A world with questions alone might not function for a day.

Now choose between two achievements, only one of which history will remember:

First, you pose a genuinely good question without answering it. Successive generations think about it for three hundred years, each becoming a little clearer-minded.

Second, you give a beautiful answer, widely quoted and taught, memorised by generations, then disproved after three centuries.

Which would you put your name beneath?

Choose and justify. First recalibrate our scales: what is a hasty answer worth, a defensible argument, an honestly calibrated I do not know? Remember the memorisers' three arguments too: bridge, map and measuring rod. They are real, not straw men.

There is no standard answer. This question never had an oval to fill. It waits on the reverse of your answer sheet for you to write.
